% Job posting: https://ca.linkedin.com/jobs/view/research-engineer-%E2%80%93-multimodal-ai-interaction-at-huawei-canada-4463665160
\documentclass[11pt]{article}
\usepackage[left=0.7in,top=0.7in,right=0.7in,bottom=0.7in]{geometry}
\usepackage{enumitem}
\usepackage[hidelinks]{hyperref}
\usepackage{titlesec}
\usepackage{array}
\usepackage{setspace}
\usepackage{needspace}
\usepackage[T1]{fontenc}
\usepackage{newtxtext,newtxmath}
\ifdefined\pdfgentounicode
  \input{glyphtounicode}
  \pdfgentounicode=1
\fi
\setstretch{1.04}
\pagenumbering{gobble}
\titleformat{\section}{\Large\scshape}{}{4em}{}[\vspace{-0.7em}\rule{\linewidth}{0.4pt}\vspace{-0.4em}]
\titlespacing*{\section}{0pt}{1em}{0.4em}
\newenvironment{cvrole}[5]{%
  \par\noindent\textbf{\large #1} | \textit{\small #2, #3}\hfill \textit{\small #4 - #5}\par
  \begin{itemize}[leftmargin=2em, itemsep=-0.3em, topsep=-0.5em]%
}{%
  \end{itemize}\par\noindent\mbox{}\par\vspace{0.1em}%
}
\newcommand{\cvName}{Nishal Stanislaus Silva}
\newcommand{\cvEmail}{nishal.silva@hotmail.com}
\newcommand{\cvPhone}{+1 613 200 2030}
\newcommand{\cvCity}{Toronto, ON}
\newcommand{\cvSite}{https://nishal.xyz}
\newcommand{\cvLinkedIn}{https://www.linkedin.com/in/nishal-silva}
\newcommand{\cvGitHub}{https://github.com/ns2max}
\newcommand{\cvStatus}{Canadian Permanent Resident, Eligible to work in Canada without sponsorship}
\newcommand{\cvTagline}{Research Scientist | Real-Time Multimodal ML, Benchmarks \& Open Datasets}
\begin{document}
\pagestyle{empty}
\begin{center}
    {\LARGE \scshape \textbf{\cvName}}{\large \textit{, PhD}}\\[1pt]
    {\small \cvTagline}\\[1pt]
    {\footnotesize\textit{\cvStatus}}\\[1pt]
    {\small
    \href{mailto:\cvEmail}{\cvEmail} \,|\, \cvPhone \,|\, \cvCity
    \,|\, \href{\cvSite}{nishal.xyz} \,|\, \href{\cvLinkedIn}{linkedin.com/in/nishal-silva}
    \,|\, \href{\cvGitHub}{github.com/ns2max}}
\end{center}

\section*{Professional Summary}
Research engineer with a PhD in real-time multimodal machine learning, most recently building and running a synchronized multi-sensor performance system that fused streaming audio, BCI/EEG biosignals (g.tec), and mixed-reality head/hand tracking across four simultaneous performers. The system was deployed in two live public concerts -- stage lighting, smoke, ten Meta Quest 3 headsets, and haptic phones -- and evaluated with a structured user study of 20 audience members and 6 performers, directly combining multimodal AI architecture with human-interaction research. This work sits on top of a broader track record in real-time embedded ML (14 ms audio-pattern inference on Raspberry Pi 4, F1 0.76, 74x faster than the DTW baseline it replaced) and multimodal sensor fusion (IMU + audio via hierarchical finite-state machines, F1 0.78), backed by rigorous benchmark design, ablation studies, and open datasets. Canadian Permanent Resident based in Toronto, available immediately, no sponsorship required.

\section*{Core Competencies}
\textbf{Multimodal AI \& Human Interaction:} real-time multi-stream sensor fusion (audio, IMU, BCI/EEG, mixed-reality head/hand tracking), synchronization of heterogeneous data streams across multiple simultaneous users, hierarchical state-machine gesture/event fusion, human-centred evaluation methodology (quantitative and qualitative user studies), Unity + Meta XR SDK prototyping, live multi-performer interactive systems. \textbf{Real-Time Systems \& Evaluation:} low-latency embedded/edge inference, sequence modelling (RNN/LSTM/GRU), digital signal processing (STFT, MFCC), frozen-protocol benchmark suite design, latency/accuracy/CPU trade-off ablations, reproducible open datasets.

\section*{Technical Stack}
Python (TensorFlow, Keras, PyTorch, scikit-learn), C++17, Unity, Meta XR SDK, Raspberry Pi / ARM CPU profiling, Elk Audio OS, Librosa, SciPy, OSC, MIDI, AWS, SQL, Docker, CI/CD, Git, LaTeX.

\section*{Education}
\textbf{PhD, Information Engineering and Computer Science} -- University of Trento, Italy (defended January 2025)\\
\textbf{MSc, Telecommunication and Electronic Engineering} -- Sheffield Hallam University, UK (2018)\\
\textbf{BEng (Hons), Electronic Engineering} -- Sheffield Hallam University, UK (2014)

\section*{Research and Work Experience}

\begin{cvrole}{Independent Researcher}{Self-directed}{Toronto, Canada}{Jan 2026}{Present}
  \item Continued publishing real-time multimodal ML research: two papers accepted for publication (Smart Drums, JAES; MIRaaS, IEEE IS2).
  \item Maintain Nebula, an open-source C++17 audio-feature-extraction library with per-feature ARM/Raspberry Pi latency benchmarks (13 GitHub stars).
  \item Use AI-assisted development workflows, including Claude Code, to accelerate research tooling and documentation.
\end{cvrole}

\begin{cvrole}{Postdoctoral Researcher}{University of Trento}{Trento, Italy}{Jan 2025}{Dec 2025}
  \item Designed and led a real-time multimodal synchronization pipeline fusing streaming audio, BCI/EEG biosignals (g.tec), and mixed-reality head/hand tracking across four simultaneous performers, in collaboration with SAE Institute Barcelona.
  \item Co-designed and ran two live multisensory concerts built on the pipeline -- synchronized stage lighting and smoke, ten Meta Quest 3 headsets, and haptic phones -- evaluated through a user study of 20 audience members and 6 performers.
  \item Owned the full ML pipeline for streaming audio and IMU/gesture data on the EU Horizon EIC Pathfinder project MUSMET: acquisition, DSP/feature engineering, training and evaluation, edge deployment, and monitoring.
  \item Published a real-time audio pattern detector running at 14 ms inference on Raspberry Pi 4, F1 0.76, 31.4\% CPU, 74x faster than the 1038 ms DTW baseline it replaced (JAES 2026).
  \item Designed frozen-protocol benchmark suites and ablation studies quantifying accuracy/latency/CPU trade-offs across model configurations.
\end{cvrole}

\begin{cvrole}{Visiting Researcher}{McGill University (IDMIL / CIRMMT)}{Montreal, Canada}{May 2024}{Aug 2024}
  \item Built a self-powered smart electric guitar combining a BNO055 IMU, HiFiBerry audio interface, and Raspberry Pi 4.
  \item Fused IMU and audio streams via a hierarchical finite-state machine, achieving F1 0.78 on a 500-event corpus (IEEE IS2 2025).
\end{cvrole}

\begin{cvrole}{Visiting Researcher}{University of Visual and Performing Arts}{Colombo, Sri Lanka}{Jan 2024}{Mar 2024}
  \item Conducted a human-centred evaluation of musicians' perception of smart musical instruments with embedded pattern detection, combining quantitative and qualitative methods (IEEE I3DA 2025).
\end{cvrole}

\begin{cvrole}{Technical Consultant}{Forestpin (Pvt) Ltd}{Colombo, Sri Lanka}{Aug 2020}{Feb 2021}
  \item Built ML and statistical pipelines for financial forensics, compliance monitoring, and risk detection on large structured enterprise data.
\end{cvrole}

\begin{cvrole}{Research Engineer}{MAS Holdings (Pvt) Ltd}{Colombo, Sri Lanka}{Jan 2015}{Jul 2020}
  \item Delivered end-to-end computer-vision and IoT systems for manufacturing at scale across a 5.5-year tenure, raising operational efficiency by up to 300\%.
  \item Built real-time IoT data-ingestion and ETL pipelines feeding C++/Python inference deployed on embedded production hardware.
  \item Introduced CI/CD and code review practices and mentored engineers across the team.
\end{cvrole}

\section*{Publications}
Silva, N., Boem, A., \& Turchet, L. (2025). Interactive IoMusT-Based Concerts: Real-Time Pattern Recognition and Audience Experience. \textit{IEEE International Conference on Immersive and 3D Audio (I3DA)}. Silva, N. \& Turchet, L. (2026). Real-Time Audio Pattern Detection for Smart Musical Instruments. \textit{J. Audio Eng. Soc.}, 74(3). Silva, N. \& Turchet, L. (2026, in press). Smart Drums. \textit{J. Audio Eng. Soc.} Silva, N. \& Turchet, L. (2026, accepted). MIRaaS. \textit{IEEE IS2}. Silva, N., Wanderley, M., \& Turchet, L. (2025). Melody and Motion. \textit{IEEE IS2}.

\section*{Service}
Peer reviewer for IEEE Access and JNSF Sri Lanka; programme committee member for IEEE IS2 and IEEE I3DA; supervised two BSc students during their thesis projects.

\section*{Highlights}
MIDI Innovation Awards 2023 Finalist ("Hot Licks", Prototypes \& Non-Commercial Software category, covered by Sound On Sound and MusicRadar); GMOA recognition (Government Medical Officers' Association, Sri Lanka, 2020) for a COVID-19 remote vitals-monitoring deployment; core contributor to the EU Horizon EIC Pathfinder project MUSMET.

\end{document}
